\documentclass[11pt,a4paper]{article}
\usepackage[margin=2.2cm]{geometry}
\usepackage{amsmath,amssymb,xcolor,enumitem,booktabs}
\usepackage[normalem]{ulem}
\usepackage[hidelinks]{hyperref}
\setlength{\parskip}{0.6em}\setlength{\parindent}{0pt}
\definecolor{newc}{HTML}{1f6f43}
\newcommand{\new}[1]{\textcolor{newc}{#1}}
\newcommand{\old}[1]{\textcolor{gray}{\sout{#1}}}
\newenvironment{greenblock}{\color{newc}}{}
\title{Proposed change, the loop, how a run is scored, and the scope}
\date{18 November 2026, version 1}
\begin{document}\sloppy
\maketitle

\section*{How to read this}
This proposal follows the margin comments on pages 29 to 35 of
\texttt{main.pdf}, Sections V-I to V-K of \texttt{3.ProblemFormulation.tex}.
Part~1 lists what is already fixed. Part~2 is what I checked in the code.
Parts~3 to 5 go through the three subsections and show only the paragraphs
that change. Black text stays, as context. \new{Green} is new and
\old{struck} text goes. Part~6 is what changes elsewhere and Part~7 the
questions. Nothing below is in the thesis yet.

Where the text and the code differ, the text follows the code (the rule
from 2026-11-18).

\section*{Part 1. Already fixed (your go-ahead in chat)}
\begin{itemize}[nosep]
  \item The broken sentence in ``The research question, formally''. A
        \texttt{\% [A55]} tag in the middle of line 1566 hid ``The tasks can
        still differ, because where a line is picked and how''. The tag now
        ends the line. The Introduction had the same fault on line 103 (``It
        is split into the wait for a robot, travel and'' was hidden), fixed
        the same way.
  \item \texttt{\textbackslash Nref} in the middle of a sentence printed
        ``The storage update (8)''. All mid-sentence uses are now
        \texttt{\textbackslash nref}, nine in the Problem Formulation and three in
        Appendices C and J. The three that start a sentence stay.
\end{itemize}

\section*{Part 2. What I checked}
\begin{itemize}[nosep]
  \item \textbf{Stall.} The code flags a stall when, in the measured window,
        no task at all is finished for longer than the stall length
        (\texttt{longest\_completion\_gap}, \texttt{evaluator.py} line 282).
        The longest time one agent stood still is logged only as a
        diagnostic. The thesis says one agent standing still, and so do
        Table~B.1, Appendix~K and the figure in Appendix~H.4.2.
  \item \textbf{Order service time.} The code averages only over orders
        released after the warm-up (\texttt{order\_metrics}). The thesis
        says ``the orders that left by $T$''.
  \item \textbf{Mean crowding.} The code averages over every step, the
        warm-up too, as the formula says. Every other measure uses only the
        steps after the warm-up.
  \item \textbf{Backlog.} The code counts waiting order \emph{lines}, not
        orders, as the notation table says. The text says orders.
  \item \textbf{Edge cost.} The Environment section names it $c(v,w)$ and
        sets it to 1 in every warehouse. Scope calls it $c(e)$ and says
        turning is inside it, which nothing does.
  \item \textbf{$\Delta$.} Appendix~G sets $\Delta=100$ as half the rotation
        period of the popular SKUs, then checks it with a sweep of 50 to 400
        steps. It was not chosen by the sweep.
  \item \textbf{$T$ and the warm-up} are never defined in the Problem
        Formulation. They first appear in ``Keeping up''.
  \item \textbf{$\hat\omega_t(v)$} is the blocked part of the service-time
        split, averaged over the picks from shelf $v$ ([A93]).
  \item \textbf{RQ4} is no longer conditional ([A102]). Scope still says
        ``if RQ4 is not reached''.
\end{itemize}

\section*{Part 3. How the two layers close the loop}

The two layers of the warehouse run at different speeds. Routing acts on
every timestep. Only every $\Delta$ timesteps\old{, at the storage epochs,}
does new stock arrive and the storage rule place it, through the storage
update (8).

The two layers influence each other in both directions, and that is what
closes the loop. Forward, storage shapes routing. Placing stock changes which
shelves a pick can take from under the stock condition, and so changes where
the agents have to go.

Backward, routing shapes storage\new{, and only congestion-aware storage has
this path}. When many picks go to one shelf, the agents wait their turn
there. \old{How long picks stand still is measured at each shelf in
$\hat\omega_t$.} \new{The shelf's standing-still estimate $\hat\omega_t(v)$
records it.} At the next storage update the rule uses it to decide where new
stock goes, and every new order line uses it to choose the shelf it picks
from. \old{This is not the crowding an agent sees around itself, the crowding
feature of Figure~10, which only the controller reads. Storage reads standing
still instead, because a crowd that keeps moving costs no time.} \new{Storage
reads standing still, not the crowding feature of Figure~10, because a crowd
that keeps moving costs no time.}

\new{Figure~11 carries the tea example of Figure~3 one storage period
further. Picks from $A$ stood still 5 steps on average and picks from $C$ 1
step. $A$ is 3 moves from the workstation and $C$ is 6. So $A$ scores
$3+5=8$ and $C$ scores $6+1=7$, and the next line takes its tea from $C$.}
\old{Figure~11 carries the tea example of Figure~3 one storage period further,
with the numbers of the pickup example in Section V-D.}

\old{With fixed storage only the forward path exists. Demand storage places
stock by the orders, and the orders come from customers, not from how the
agents move. So with demand storage the coupling still runs one way, as a
gain that depends on the controller.}
\new{Fixed storage never moves where stock goes, and demand storage follows
the orders. The orders come from customers, not from how the agents move. So
under both rules the coupling runs one way, from storage to routing.} The
backward path is what congestion-aware storage adds, because only that rule
reads the time agents stand still. \old{The loop is closed only in RQ2, which
tests it at the main load and at a higher load where the aisles are
crowded.}

Separating the two clocks has one consequence that constrains the study.
Between two storage updates the routing traffic settles into a steady
pattern, and the storage rule only sees \old{its average per step over the
period just ended} \new{what happened in the recent periods}. If $\Delta$ is
too small, the traffic has not settled before the next update, so
\old{$F$ reacts to noise} \new{the storage rule reacts to chance}\old{, and
the shelves the agents pick from keep changing while they are still
learning}. \new{The shelves the agents pick from then keep changing, which
also makes routing harder to learn.} If $\Delta$ is too large, the numbers
$\hat\rho$ and $\hat\omega$ are out of date by the time they are used\new{,
and new stock arrives late}. \old{The choice of $\Delta$ is therefore a
parameter of the study, set by a sweep (Table B.1), not a constant fixed in
advance.} \new{$\Delta$ is therefore set from how fast the popular SKUs
change, and a sweep of periods checks it (Appendix~G).}

\section*{Part 4. How a run is scored}

A run is scored on the whole warehouse, not on each layer. (Unchanged.)

The score follows lifelong MAPD. H. Ma et al. (2017) judge a lifelong MAPD
algorithm by its service time, and they count an instance as solved only if
the service time stays bounded. This thesis does the same. \old{First we
check that the fleet keeps up with the arriving tasks. Then the result is the
mean time an order takes, from its release until it leaves at a dock. Three
things explain it: the time of single tasks, how that time splits into
waiting, moving and being blocked, and how crowded it was. Put-aways and the
time the controller needs to decide are reported next to it.}

\begin{greenblock}
A run lasts $T$ timesteps, the horizon. It starts with no orders waiting, so
its first steps do not look like a warehouse that has been running for a
while. These steps are the warm-up. The steps after it are the measured
window, and every measure below counts only those, unless it says otherwise.
\end{greenblock}

\textit{Keeping up.} A fleet keeps up when it finishes tasks about as fast as
they are released. If it falls behind, unfinished tasks pile up, and an
average over the finished ones no longer describes the run. This is the
bounded service time that H. Ma et al. require of a solved instance.
\old{The queue is the backlog $B_t = \ldots$, the tasks released but not yet
finished and the orders still waiting for stock.}
\new{The unfinished work is the backlog. It holds the tasks released but not
yet finished, and the order lines still waiting for stock,
$\mathcal{Z}^{\mathrm{wait}}_t$. Formally,
$B_t = |\mathcal{Q}_t| + |\mathcal{B}_t| + |\mathcal{Z}^{\mathrm{wait}}_t|$.}

\old{Throughput is the number of tasks finished per timestep. It is the
measure lifelong MAPF uses for warehouses (J. Li et al., 2021) and the pick
rate that Krnjaic et al. (2022) optimise. The check compares it with the rate
at which tasks were actually released in the run, not the expected rate,
because orders arrive at random and one seed releases more than another.}
\new{The keep-up check compares the tasks the fleet finished with the tasks
actually released in the run, not with the expected number, because orders
arrive at random and one seed releases more than another.}

The check counts every task the fleet finishes in the measured steps, also
one released before them. A fleet that fell behind earlier still works
through its older tasks, and that work is what shows whether the
\old{queue}\new{backlog} grows or holds. So the measure is the share of the
arriving work the fleet serves. With 900 tasks released in the measured steps
and 600 tasks finished in them, the share is two thirds and the
\old{queue}\new{backlog} grows by 300.

A pick counts as released when its order arrives, also while its stock is not
yet on a shelf. The order is work the fleet owes, and the backlog already
counts it. Left out, a fleet whose put-aways fail would release almost no
picks and seem to keep up. Such a fleet with 800 \old{orders}\new{order
lines} waiting for stock and 10 picks finished serves about one pick in
eighty.

Picks and put-aways are checked separately\ldots{} (Unchanged.)

Formally, write $\mathcal{R}^{\mathrm{pick}}_T$ for the picks released in the
measured window\old{, with one for each order line from the window that still
waits for stock at the horizon, and}\new{. An order line from the window that
still waits for stock at $T$ counts as one more released pick. Write}
$\mathcal{D}^{\mathrm{pick}}_T$ for the picks finished in it, whenever they
were released, and the same for put-aways. The served share of a kind is the
second over the first\old{, and the}\new{. The} run passes the keep-up check
if
\[
  |\mathcal{D}^{\mathrm{pick}}_T| / |\mathcal{R}^{\mathrm{pick}}_T| \ge f_{\mathrm{up}}
  \quad\text{and}\quad
  |\mathcal{D}^{\mathrm{put}}_T| / |\mathcal{R}^{\mathrm{put}}_T| \ge f_{\mathrm{up}} , \qquad (27)
\]
where the threshold $f_{\mathrm{up}} \le 1$ leaves room for tasks still under
way at the horizon. \old{The throughput $\Lambda_T$ counts the same finished
tasks, picks and put-aways together, per step after the warm-up. A fleet that
fell behind is credited with the older tasks it finishes.} The backlog at
the horizon splits the same way into $B^{\mathrm{pick}}_T$ and
$B^{\mathrm{put}}_T$, with the \old{orders}\new{order lines} still waiting
for stock in $B^{\mathrm{pick}}_T$, and is reported next to the check.
\old{A run that does not keep up is reported by its served share and is not
compared on service time.} The check looks at one finite horizon, so passing
it does not prove that the backlog stays bounded. A run can finish most of
its tasks while the backlog still grows slowly. The backlog over the run is
therefore shown for the conditions in the main result, so a growing backlog
is visible.

\begin{greenblock}
The throughput $\Lambda_T$ is the number of tasks finished per step in the
measured window, picks and put-aways together. It is the measure lifelong
MAPF uses for warehouses (J. Li et al., 2021) and the pick rate that Krnjaic
et al. (2022) optimise. It counts the same finished tasks as the check, so a
fleet that fell behind is credited with the older tasks it finishes. A run
that does not keep up is reported by its served share and its throughput,
and is not compared on service time.

The number of orders finished per step is reported next to throughput.
Throughput counts tasks, and a storage rule that splits more lines over
several shelves finishes more tasks without serving more customers.
\end{greenblock}

A fleet can also jam for a stretch and then catch up. Its throughput
recovers, so it passes the check, and its service time carries the jam. The
jam is counted on its own as a stall, a stretch \old{in which an agent with a
task stands still longer than a set length}\new{of the measured window in
which no agent finishes a task for longer than a set length} (Appendix
H.4.2). \old{It is counted because it says something about the controller
rather than about the load.} \new{A healthy fleet finishes tasks all the
time, so a long gap points to the controller rather than to the load.}

\textit{Service time.} The result is the mean order service time\ldots{}
(Unchanged.)

\old{Formally, an order $z$ is released at $t_z^{\mathrm{released}}$, the
release time all its lines share, and leaves when its shipment is delivered,
at $t_z^{\mathrm{finish}}$.} \new{Formally, write $t_z^{\mathrm{released}}$
for the step order $z$ comes in, shared by all its lines, and
$t_z^{\mathrm{finish}}$ for the step its shipment is delivered.} Its service
time is $\zeta^{\mathrm{ord}}_z = t_z^{\mathrm{finish}} -
t_z^{\mathrm{released}}$ (28), and with $\mathcal{Z}^{\mathrm{done}}_T$ the
orders \new{released in the measured window} that left by $T$, the result is
the mean (29).

\textit{Where the time goes.} (First paragraph unchanged.)

Formally, let $\omega_j(t) \in \{0,1\}$ be 1 when the agent serving $\tau_j$
is blocked at $t$, and let $\zeta^{\mathrm{travel}}_j$ be the number of
timesteps in which it moves. Then (30)\old{,}\new{.} The split is taken over
the picks, the trips that fetch an order's lines, because that is where
storage acts. Averaging each part over the \old{completed picks}\new{picks
completed in the measured window, $\mathcal{C}^{\mathrm{pick}}_T$,} gives
three means that add up to the mean service time of picks. The blocked part
averaged in this way is the mean blocked time $W_T$ (31). \new{The
standing-still estimate $\hat\omega_t(v)$ of a shelf is built from the same
blocked steps, averaged over the picks from that shelf.}

Each part points to one cause. The wait before assignment grows when the
fleet is too small for the load. It also holds any time an order waits for
stock, which is reported on its own when it is not small. Travel shrinks when
the storage rule puts popular items close to \old{where they are
needed}\new{their workstations}\old{, and grows when it puts them far from
the docks}, so this is where demand and congestion-aware storage should
show. Blocked time grows with congestion. The split is this thesis's own.

\old{\textit{Orders finished.} The number of orders finished per timestep is
reported next to throughput\ldots{}} (Moved up, after throughput.)

\textit{Crowding.} Congestion is measured by how crowded the space around the
agents was. \old{Each agent already carries this reading, the congestion
feature (21), the share of its field of view occupied by other agents.}
\new{The measure is the crowding feature (21), the share of an agent's field
of view taken by other agents. It is computed for every agent under every
controller, also where the agent does not see it.} Averaging it over every
agent and every timestep gives one number for the run, the mean crowding
(32). In Figure~10 the three agents read 0, 0.25 and 0.25, so that timestep
contributes about 0.17. \new{The reading follows the crowd density of CRAMP
(Pham and Bera, 2024), here averaged over the run (Appendix C.3).}

\old{The reading follows the crowd density of CRAMP (Pham and Bera, 2024),
averaged into a reported outcome (Appendix C.3).}

Crowding and blocked time are read together. Crowding without blocking costs
nothing, as with a line of agents that keeps moving. Crowding with blocking
is congestion that costs time. \old{The risk that feedback funnels traffic
onto one aisle is not scored on its own. Congestion-aware storage guards
against it, because it avoids shelves where picks stand still, and when it
happens it shows up as crowding and blocked time.} \new{Congestion-aware
storage could send so much stock to a quiet aisle that the quiet aisle
becomes the crowded one. No measure scores this on its own. It shows up as
crowding and blocked time there, and the next storage update then moves new
stock away again.}

\textit{Reported measures.} Two \old{measures}\new{costs} are reported next
to the result. (Rest unchanged.)

Table~VI sums up the score. (Text unchanged.) The table changes as follows.

\begin{center}\small
\begin{tabular}{lll}
\toprule
Quantity & Role & Source \\
\midrule
served share (27) & check & H. Ma et al. (2017), J. Li et al. (2021) \\
\new{backlog $B_T$} & \new{check} & \new{H. Ma et al. (2017)} \\
\new{stalls} & \new{check} & \new{this thesis} \\
served share and $\Lambda_T$ & result when a run does not keep up & Krnjaic et al. (2022) \\
\new{orders finished per step} & \new{explains $\Lambda_T$} & \new{this thesis} \\
$\bar\zeta_T$ (29) & result & after Krnjaic et al. \\
mean service time of picks, $\zeta_j$ (13) & explains the result & H. Ma et al. (2017) \\
three parts of $\zeta_j$, $W_T$ (30), (31) & explains the result & \old{queue split}\new{service-time split}, this thesis \\
$\bar\delta_T$ (32) & congestion & after Pham and Bera (2024) \\
mean service time of put-aways & reported cost & H. Ma et al. (2017) \\
$\kappa_T$ (33) & reported cost & H. Ma et al. (2017) \\
\bottomrule
\end{tabular}
\end{center}

\textit{The decision problem.} (First paragraph unchanged.)

The thesis does not search for the best pair. \new{The problem is stated so
that the comparison has a clear target.} It compares three rules under three
controllers, and the gain test below asks whether a rule beats fixed storage
under one controller.

\textit{The research question, formally.} The thesis compares fixed storage
with demand and congestion-aware storage under each controller, and the
comparison is paired. Under one seed, \old{as Section V-D, Orders and tasks,
says,} fixed storage and the rule $F$ receive the same customer orders\new{,
because the orders never depend on the stock}. The tasks can still differ,
because where a line is picked and how many tasks it splits into depend on
where the stock is. (Rest of the paragraph and (35), (36) unchanged.)

This gain test is a paired $t$-test over seeds\new{, paired because both runs
of a seed meet the same orders}. How the gain holds as the load grows is RQ3.

Controllers are compared on the gain as a percentage of their mean service
time under fixed storage, the form in which storage studies report a gain
(Puiseau et al., 2022)\new{, for the reason given with the research
questions}. \old{A controller whose service is slow under fixed storage has
more to save, so the gain in timesteps alone would favour the controller that
routes worst.} Formally, RQ1 compares (37) between controllers, with its 95\%
interval over seeds as in the gain test (36).

The study then asks how the gain depends on three factors. The first is the
controller, how many agents one decision of $\pi$ covers (Table IV). The gain
is read together with its split\old{, because a fleet that runs close to its
limit turns the same saving in travel into a larger cut in the wait before
assignment}. \new{A fleet that is nearly always busy has tasks waiting for an
agent. Shorter trips free agents sooner, so those tasks wait less too. The
same saving in travel then shows up as a larger cut in the wait before
assignment.} The second is whether the storage rule is congestion-aware,
congestion-aware storage against demand storage. The third is the load, set
by the fleet size $m$ and the expected number of order lines released per
timestep $\lambda_{\mathrm{task}}$. On the same sweep, throughput per agent,
$\Lambda_T/m$, shows whether each agent finishes fewer tasks as the fleet
grows, and whether that loss is faster when one decision covers fewer agents.
For RQ4, the decentralised controller is also run with communication over
the communication graph (22) and without it. Only that controller has
messages, so this comparison \old{sits outside the main factorial design}\new{is
made apart from the comparison of rules and controllers}.

\section*{Part 5. Scope and assumptions}

\begin{itemize}[nosep]
  \item \textbf{No physical model.} There is no model of load, battery or how
        an agent drives. \old{Turning and speeding up are counted inside the
        edge cost $c(e)$.} \new{Turning and speeding up are not modelled.
        Every move takes one timestep.}
\end{itemize}

Tasks carry no deadline. A late delivery only lengthens service time, and no
single task can be said to have failed. This is narrower than
time-constrained formulations such as Knippenberg et al. (2021), where a
missed deadline is infeasible or penalised. \old{Adding deadlines would put a
deadline in every task, and either forbid a late task or count late tasks
next to the result.} \new{Adding deadlines would let a late task fail.}

Assignment is also permanent once made\ldots{} Either change alters what
counts as a solution, so it would change this chapter and not the method.
(Unchanged, see Question~5.)

Communication is built in and can be switched on or off, but RQ1 to RQ3 do
not depend on it. \old{If RQ4 (Section II-B) is not reached, the results say
nothing about communication.} \new{Only RQ4 uses it.}

Each item above limits what the results can claim. (See Question~6.)

\section*{Part 6. What changes elsewhere}
\begin{itemize}[nosep]
  \item \texttt{styles/eqnames.tex}: \texttt{eq:congestion} is named
        ``crowding feature'', not ``congestion feature'', to match the text
        ([A93] renamed it).
  \item \texttt{A.Notation.tex} and \texttt{2.Notation.tex}: a row for the
        warm-up is not needed (words only). A row for
        $\mathcal{C}^{\mathrm{pick}}_T$, ``picks completed in the measured
        window''. The $T$ row already says ``the horizon''.
  \item The stall: \texttt{B.Parameters.tex} (row ``stall''),
        \texttt{K.Checks.tex} line 32 and the figure and calibration in
        Appendix~H.4.2 say one agent standing still. They follow the code
        once you choose (Question~1).
  \item \texttt{B.Parameters.tex}, the $f_{\mathrm{up}}$ row still says
        ``share of the release rate the throughput must reach''. It becomes
        ``share of the released work the fleet must serve''.
  \item Figure~11 (\texttt{shelf\_loop.png}): panel~(a) also shows where
        demand storage puts the new tea and where the next line picks (both
        at $A$), and ``6 edges'' becomes ``6 moves''. Redrawn from the
        script (Question~3).
  \item \texttt{GAPS.md}: one new entry (the next free number, [A107] or
        [A108] depending on the observation review), with \texttt{\%} tags in
        the chapter.
\end{itemize}

\section*{Part 7. Questions}
\begin{enumerate}
  \item \textbf{Stall.} The code counts a stall when the whole fleet finishes
        nothing for too long. The text, Table~B.1 and Appendix~H count one
        agent standing still. I wrote the code's version. Keep that, or should
        the code change?
  \item \textbf{Crowding window.} The mean crowding counts the warm-up, every
        other measure does not. Keep it (the formula says so), or move it to
        the measured window in the code?
  \item \textbf{Figure~11.} Why does the new tea go to the middle shelf and
        not to $C$? It depends on the shelf score, which sits in
        Implementation. Should the caption say ``the quiet aisle'' only, or
        should the figure show the score?
  \item \textbf{Order of the loop section.} The point of the section is that
        only congestion-aware storage closes the loop. Should that sentence
        open the section, before forward and backward?
  \item \textbf{Permanent assignment.} The text says what Ma et al.\ allow,
        not why we chose otherwise. Is the reason that it keeps the wait for
        an agent apart from travel, or something else?
  \item \textbf{Last sentence of Scope.} ``Results about decentralised
        control assume that every allowed move succeeds.'' Every controller
        assumes this. Did you mean decentralised control suffers most when
        moves can fail? If so, say why, or drop ``decentralised''.
  \item \textbf{Throughput per agent.} $\Lambda_T/m$ came from the commons
        question, folded into RQ4 by [A33]. RQ4 is now communication. Which
        question does it serve, RQ3, or should it go?
  \item \textbf{Token Passing bullet.} ``One agent per workstation at a time''
        is about one router. Keep it in Scope, or move it to the Method,
        where Token Passing is described?
  \item \textbf{Decision problem.} $\mathcal{F}$ is every storage rule, but
        $\Pi$ is only the controllers studied. Should $\Pi$ also be every
        controller, so the problem is stated apart from the method?
  \item \textbf{$W_T$ next to the workstation $W_1$.} Keep, or rename the mean
        blocked time? Check \texttt{notation\_renames.md} first.
\end{enumerate}

\end{document}
